\chapter*{Kurzfassung}
\addcontentsline{toc}{chapter}{Kurzfassung}

Diese Arbeit behandelt die schwach überwachte Mutationsvorhersage aus
Gigapixel-H\&E-Bildern (mit Hämatoxylin und Eosin gefärbt) von
Lungenadenokarzinomen (LUAD, dem häufigsten Subtyp des
Lungenkarzinoms): einem einzigen binären Label pro Gen werden etwa
$10^9$ Pixel zugeordnet, ohne Annotation auf Subregion-Ebene.
Bestehende Systeme benötigen Zehntausende proprietärer Schnitte und
liefern keine Erklärung dafür, \emph{warum} eine Mutation
vorhergesagt wird --- ein zentrales Hindernis für klinisches
Vertrauen. Das zugrunde liegende Problem ist seinerseits
\emph{stochastisch, mehrdeutig (viele-zu-viele) und schwach
gelabelt}: dieselbe Morphologie kann mehreren Genotypen entsprechen
und umgekehrt, weshalb die Aussagekraft eines Modells pro Gen und
pro Schnitt belegt werden muss, nicht aus einer einzigen
Kohortenmetrik.

Die zentrale Hypothese lautet, dass \emph{strukturiertes
Domänenwissen} --- konkret die sechs durch Experten definierten
histologischen Wachstumsmuster des LUAD
(lepidic, acinar, papillary, micropapillary, solid, cribriform),
kodiert mittels der Theorie unscharfer Mengen --- in eine
Deep-Learning-Pipeline integriert werden kann, um gleichzeitig
wettbewerbsfähige Mutationsvorhersagen unter Datenknappheit
($<\!700$ öffentliche Schnitte) zu erzielen und die Vorhersagen
\emph{interpretierbar} zu machen, mittels eines mehrstufigen,
klinisch verifizierbaren Inspektionsrahmens --- und darüber hinaus
\emph{erklärbar} durch wissensgraph-basierte
natürlichsprachliche Begründungen --- in einem Ausmass, das kein
End-to-End-Black-Box-Ansatz bietet; das Urteil fällt
\emph{gen-konditional} aus.
Nach den Prinzipien der Design Science Research werden drei
Artefakte und ein integratives System (\lchai{},
\emph{Lung Cancer Histologic Analysis with AI}) entworfen und
evaluiert.

\textbf{Artefakt~1 (FuzzyArcLoss~V2)} ist eine probenweise
unscharfe Modulation des ArcFace-Winkelmargins, welche die
graduellen morphologischen Übergänge zwischen Musterklassen
respektiert; im Vergleich von 18 Verlustfunktionen erreicht es
Platz~1 (Macro-F1 $= 92{,}31\,\%\!\pm\!2{,}04$, $p\!=\!0{,}001$,
Cohen's~$d\!=\!1{,}00$). Sein Vorteil ist
\emph{abhängig von morphologischen Klassengrenzen, nicht von der
Klassengrösse}: der Gewinn manifestiert sich auf visuell
mehrdeutigen, benachbarten Klassen unabhängig von deren Grösse ---
eine Übertragungsregel, die sich auf Gleason-Grading,
Pap-Zytologie und Dermatopathologie verallgemeinern lässt.

\textbf{Artefakt~2 (PI-ABMIL)} verkettet einen 512-dimensionalen
CTransPath-Embedding-Vektor mit den 6-dimensionalen unscharfen
Musterwahrscheinlichkeiten aus Artefakt~1 vor der
Gated-Attention-Aggregation und verankert damit die
Schnitt-Ebene-Argumentation in klinisch validierten Mustern.

\textbf{Artefakt~3 (FC-MIL)} ersetzt die additive Attention durch
ein erlerntes 2-additives Choquet-Integral, dessen 21~Parameter
(sechs Shapley-Werte, fünfzehn paarweise Interaktionen)
supraadditive Wechselwirkungen zwischen morphologischen Kategorien
erfassen, welche Standard-Attention nicht repräsentieren kann.

Die Evaluation auf TCGA-LUAD (\emph{The Cancer Genome Atlas
Lung Adenocarcinoma}; 687~Schnitte, 668~Patientinnen und Patienten,
5-fache patientstratifizierte Kreuzvalidierung) über sechs
Treibermutationen (\textit{TP53}, \textit{EGFR}, \textit{KRAS},
\textit{STK11}, \textit{KEAP1}, \textit{RBM10}; absteigende
Prävalenz) zeigt, dass die Nützlichkeit des Musterwissens
\emph{gen- und aggregator-konditional} ist und sich in drei
strukturelle Regime zerlegt:

\textbf{Regime~(i) --- Encoder-Sättigung.}
Der vorgelernte Encoder erfasst bereits die relevante Morphologie
und der Musterkanal fügt nichts hinzu (\textit{TP53}, \textit{EGFR}).

\textbf{Regime~(ii) --- Inter-Muster-Kookkurrenz.}
Das diskriminative Signal lebt in paarweisen Musterinteraktionen
und nicht in einem einzelnen Muster; nur ein Aggregator, der
gemeinsame Aktivierungen liest, erfasst es (\textit{KRAS}, durch
das Choquet-Integral von FC-MIL).

\textbf{Regime~(iii) --- ausserhalb der Ontologie.}
Keine Repräsentation innerhalb des Sechs-Muster-Simplex
(ANORAK~\cite{anorak2024}) trägt das Signal; das Regime zerlegt
sich in drei benannte Mechanismen --- \emph{label-space
mismatch}~\cite{lipton2018labelshift}, \emph{no morphological
projection} und
\emph{collapse-to-prior}~\cite{bowman2016posteriorcollapse} ---,
welche zusammen begrenzen, was eine reine H\&E-basierte
Mutationsvorhersage erreichen kann (\textit{STK11}, \textit{KEAP1},
\textit{RBM10}).

Um diese Grenzen sichtbar zu machen, instrumentiert \lchai{} einen
sechsstufigen Interpretierbarkeits-Stack (Attention,
Mustertüberlagerungen, SHAP, Choquet-Shapley-Werte, Kanal-Ablation,
unscharfe linguistische Labels) sowie einen LLM-Erklärungsagenten,
verankert in einem kuratierten Wissensgraphen aus den Ontologien
NCIt und MONDO: unscharfe Labels regeln, \emph{wie} das LLM spricht,
der Graph regelt, \emph{was} es behaupten darf --- ein doppelter
Anti-Halluzinations-Mechanismus.
Eine ergänzende kohortenbezogene Conclusive\,/\,Inconclusive-Markierung
(aus AUROC abgeleitet) steuert den schnitt-ebenen Erklärungspfad
und trennt populationsbezogene Signal-Existenz von der
schnitt-spezifischen morphologischen Konfidenz.

Die Hypothese wird \emph{teilweise gestützt}: Musterwissen liefert
eine Interpretierbarkeit, die kein bisheriges
Mutationsvorhersagesystem bietet, sowie einen echten
inter-musterlichen Vorhersagegewinn bei \textit{KRAS}; bei den
übrigen Genen liegt die Schranke in der 6-Muster-\emph{Eingabe-%
repräsentation} selbst (dem ANORAK-Label-Raum), nicht im
nachgelagerten Modell (Verlust\,/\,Merkmalsfusion\,/\,Aggregator).
Die Hauptansprüche dieser Arbeit sind demzufolge nicht klinische
Genauigkeit --- dafür existieren bereits foundation-skalige
Kohorten von 60\,000+ Schnitten ---, sondern die folgenden zwei,
destilliert aus den sechs feingranularen Beiträgen C1--C6 in
Kapitel~\ref{ch:introduction}:

\textbf{Anspruch~(i) --- methodisch.}
Unscharfe Logik, kohärent angewandt über Verlust, Repräsentation
und Aggregation, macht eine auf öffentlichen Daten basierende
Pipeline transparent durch einen wissensgraph-gestützten
sechsstufigen Interpretierbarkeits-Stack.

\textbf{Anspruch~(ii) --- empirisch/analytisch.}
Die gen-spezifische Taxonomie aus den schnitt-bezogenen Ablationen
auf TCGA-LUAD begrenzt, was ein rein morphologie-basiertes Modell
aus H\&E allein erreichen kann, und identifiziert jene Gene, für
welche die Fusion zusätzlicher Datenquellen (z.\,B.\
Genexpressionsprofile oder Radiologie), nicht weitere
architektonische Verfeinerung, der natürliche nächste Schritt ist.

Zusammen positionieren diese Ansprüche den vorgeschlagenen Ansatz
als komplementäres Paradigma zur Brute-Force der
Foundation-Modelle --- eines, in dem eine Pathologin oder ein
Pathologe das Reasoning des Modells verstehen, verifizieren und
\emph{korrigieren} kann, und in dem Interpretierbarkeit gerade in
jenen Fällen am wertvollsten ist, in denen der Prädiktor irrt.

\paragraph{Schlüsselwörter.}
Unscharfe Logik, Winkelmargin-Verlust, Multiple-Instance-Lernen,
Choquet-Integral, Shapley-Werte, histopathologische Muster,
Lungenadenokarzinom, Mutationsvorhersage,
Genotyp-Phänotyp-Lücke, Design Science Research,
Computational Pathology, erklärbare KI.
